\input{../../../assets/programming-editorial.tex}
\renewcommand{\studiumcoverphoto}{../../../assets/covers/vol04.jpg}
\urldef{\studiumsourceurl}\url{https://commons.wikimedia.org/wiki/File:286_AT_motherboard.jpg}
\urldef{\studiumlicenseurl}\url{https://creativecommons.org/licenses/by-sa/4.0}
\newcommand{\studiumcustomtitle}{
\studiumcover{C y C++}{Memoria, recursos y programación de sistemas}{VOLUMEN 4}
\chapter*{Créditos y alcance de esta revisión}
Daniel Miguel Tejedor. Biblioteca de Programación, Studium. 2026.

Esta copia conserva los 75 capítulos de la colección original. La revisión de fuentes, cobertura y soluciones sigue abierta. Un cálculo aritmético reproducido no verifica un programa. Esta edición no cuenta con aprobación académica independiente.

\noindent\textbf{Fotografía de portada.} Mister rf / Wikimedia Commons / CC BY-SA 4.0. 286 AT motherboard.jpg.\par\studiumsourceurl

\noindent\textbf{Licencia y transformaciones.} CC BY-SA 4.0.\par\studiumlicenseurl

Se conservan los bytes originales. Colocación proporcional y recorte en la portada. La imagen y su recorte conservan la licencia indicada.

\noindent\textbf{Cómo estudiar.} Escribe los ejemplos en archivos, ejecuta las pruebas indicadas y compara la salida observada con la esperada. Consulta las soluciones después de intentarlo. Los bloques Terminal contienen comandos; Output muestra resultados y Diagnostic contiene errores.\par\clearpage
}
